Magnetic actuation applies remote forces and torques to a flexible body through a field and its spatial gradient \citep{abbott,kim}. A useful control model must carry these inputs through mechanical reduction, identification and numerical implementation without changing the dynamical object at each step. This paper studies that chain for a planar clamped--free magnetic catheter model. Negligible gravity motivates the application setting; small slope, no contact and no surrounding-fluid load define the present benchmark. More general rod mechanics remain the starting point for extending these assumptions \citep{simo,antman}.

The nonlinear target retains the complete coupled cosine magnetic potential. Its higher-order terms are not canceled to produce independent linear oscillators: they are the source of the nonlinear frequency variation used in the persistence analysis. The controller compensates an identified conservative residual and proportional damping around that nominal target. The central question is therefore whether the same learned compensation can be connected to a dimensionally correct difference scheme, a physical magnetic command and a precise statement about the resulting discrete dynamics.

\subsection{Three linked subproblems}
\label{fw:three-questions}
\textbf{P1: derive and implement the difference scheme.} Starting from a compatible Birkhoff representation, which endpoint variables admit a generating function, and which matrix update follows? The mechanical mass matrix, phase-flow Jacobian and endpoint transformation have different dimensions. Their distinction matters when a scalar-looking formula is implemented for several modes. Proportional damping also distinguishes a weighted structure-preserving map from ordinary midpoint applied to the damped vector field. We make these constructions explicit and test the noncommuting matrix operations before using the nonlinear catheter model.

\textbf{P2: identify and physically realize residual compensation.} Can state observations determine a conservative residual and damping estimate that preserve the nominal coupled magnetic potential? The identification, controller and simulated plant must use the same coordinates, while the learned controller must not read the unknown coefficients used to generate synthetic data. A second requirement is physical: two generalized modal commands must pass through a common force--torque interface and then a local field--gradient reconstruction. A small identification residual does not by itself verify either sufficient excitation or hardware feasibility.

\textbf{P3: establish the structure and convergence of the controlled scheme.} Once feedback is placed inside the same midpoint equation, which vector field does the scheme approximate, what structure survives identification error, and how does the nonlinear solve affect convergence? The finite-time error must be measured against the same actual closed loop. Its deviation from the nominal target is a separate modeling error. Exact symplecticity is available on the autonomous conservative branch; discrete KAM additionally requires analytic domains, nonresonance, twist and a quantitative smallness condition. These implications cannot be inferred from a phase portrait or a small energy drift.

\subsection{One physical and statistical model throughout}
The benchmark is built from a literature material reference, manufacturer magnet specifications and declared geometric and damping choices. The Sylgard 184 reference is associated with its published formulation and curing conditions \citep{cho2021pdms}; it is not a measurement on the present beam assembly. The tip core contributes translational and rotary inertia before the biased eigenproblem is solved. Those inputs determine the modal scales, physical decoder and magnetic input map used in all three subproblems. Their provenance is given in Section~\ref{sec:derivation} and Supplement~\ref{sec:supp-mechanics}.

The mechanical chain is kept distinct from the learning hypothesis class. Hamiltonian neural networks learn a scalar energy whose derivatives determine a vector field \citep{greydanus}; separate dissipation and actuation terms retain a meaningful power balance \citep{zhong,desai}. Complex coordinates organize the real position and momentum channels, and complex features can supply structured representations \citep{trabelsi}. The real energy depends on a complex coordinate and its conjugate. The analytic continuation required by KAM instead treats the underlying complexified coordinates as independent variables. Neither this notation nor a universal approximation theorem supplies a validated derivative error for an arbitrary fitted network \citep{hornik,zhang}.

The implemented experiment fixes the known kinetic energy and learns an analytic potential residual together with a scalar damping coefficient from integral state observations. It uses a single declared magnetic feature, allowing identification error, numerical error and input reconstruction to be checked separately. A larger analytic readout remains part of the theoretical framework, but the reported computations do not claim unrestricted hidden-layer training or a new complex-versus-real network comparison. The complete nominal cosine is retained even in the single-feature test.

\subsection{Geometric analysis and the control objective}
Generating functions provide established symplectic methods when their geometric and local invertibility assumptions hold \citep{feng,fengqin,hairer}. We use a constructive proportional-damping Birkhoff representation, derive its endpoint projection and time-midpoint update, and recover the ordinary Hamiltonian midpoint map in the conservative limit. The fourth-order central generating function is retained with its complete correction. These classical constructions are applied to the declared mechanical and feedback model; a new general-purpose integration method is not claimed.

The controller estimates damping and a potential correction while retaining $h_0$ as the nominal target. Its actual conservative closed-loop part is $h_c=h_0+\delta R$, where $\delta R$ is the residual identification error. Remaining damping must be retained as a nonconservative term. Consequently, time-discretization convergence can hold when exact symplecticity and classical KAM do not. Section~\ref{sec:analysis} states this separation as a controlled numerical-analysis result, including the accumulated effect of the implicit-solve residual.

The magnetic quartic normal form provides a computable twist for a conditional persistence argument. Learning-to-KAM results require hypotheses beyond a small empirical loss \citep{chen2022}; we therefore retain a deterministic coefficient and analytic-error budget. Continuous KAM offers background for the conservative target \citep{poschel}, while the principal discrete conclusion uses the generating-map framework of \citet{shang}. A finite-amplitude a posteriori certificate would additionally require a specific embedding, invariance defect and geometric bounds \citep{haro2019aposteriori}. Such a new certificate is not supplied by the numerical diagnostics in this study.

Related magnetic-robot studies couple reduced models, visual feedback and simulation to physical control \citep{francescon2025shape,tong2025navigation}. The present contribution isolates the analytical and computational connections needed before those additional sensing and hardware layers are introduced. Learning a conservative reference and tracking it is another valid objective, but it produces a different error system. That related route is retained in Supplement~\ref{sec:supp-tracking}; the main article follows the residual-compensation objective throughout.

\subsection{Contributions and organization}
The first contribution is a complete path from the Birkhoff one-form and endpoint projection to a matrix implementation, with the physical, modal and phase-space dimensions stated explicitly. The second is a shared integral-identification and magnetic-realization construction for residual compensation on the coupled nominal model. The third is an analysis of the controlled discretization that separates finite-time convergence, exact or weighted geometric structure, and conditional discrete KAM persistence. The numerical study is organized to test the corresponding interfaces rather than to combine unrelated accuracy metrics.

\begin{table*}[t]
\centering
\caption{The three subproblems within the six-section article. The same mechanical baseline and compensation objective are retained across the rows.}
\label{tab:roadmap}
\begin{tabularx}{\textwidth}{p{0.15\textwidth}YY}
\toprule
Subproblem & Construction in the article & Verification and limits\\
\midrule
P1: discrete scheme & Section~\ref{sec:derivation}: compatible Birkhoff form, endpoint projection, matrix midpoint and weighted map & Section~\ref{sec:analysis}: independent matrix-exponential references, noncommuting elimination and structure diagnostics\\
P2: learned control & Section~\ref{sec:derivation}: joint integral identification and residual compensation; Section~\ref{sec:application}: physical magnetic ports & Section~\ref{sec:analysis}: four control ablations on one plant; noiseless restricted-feature evidence\\
P3: controlled analysis & Section~\ref{sec:analysis}: feedback inside the implicit scheme and solver-residual convergence; Section~\ref{sec:kam}: discrete persistence conditions & Fixed-model and step-refinement diagnostics; the full analytic KAM budget and device calibration remain separate requirements\\
\bottomrule
\end{tabularx}
\end{table*}

Section~\ref{sec:derivation} supplies the common mechanics and the discrete and learning constructions. Section~\ref{sec:analysis} establishes the controlled scheme's finite-time properties and presents the three-part verification sequence. Section~\ref{sec:kam} connects the conservative branch to discrete KAM. Section~\ref{sec:application} realizes the same command through magnetic ports and states the remaining full-beam and clearance requirements. Section~\ref{sec:conclusion} summarizes the answers to P1--P3. The Supplement provides the detailed proofs, numerical protocols and the related tracking extension.
